\documentclass[10pt,a4paper]{amsart}
\usepackage[margin=19mm]{geometry}
\usepackage{amsmath,amssymb,mathtools}
\usepackage[scr=rsfs,cal=euler]{mathalfa}
\usepackage{xfrac}
\usepackage[unicode,hidelinks,bookmarks=false]{hyperref}
\newcommand{\ie}{i.e.\ }
\newcommand{\cf}{cf.\ }
\newcommand{\resp}{respectively\ }
\newcommand{\Subsection}{\subsection}
\providecommand{\nobreakdash}{\nobreak\hbox{-}}
\setlength{\emergencystretch}{2em}
\multlinegap0pt
\usepackage{bm}
\usepackage{bbm}
\usepackage{dsfont}
\usepackage{xspace}
\usepackage{cancel}
\usepackage{mathtools}
\usepackage{amsthm}
\makeatletter
\@ifundefined{theorem}{\newtheorem{theorem}{Theorem}}{}
\@ifundefined{proposition}{\newtheorem{proposition}[theorem]{Proposition}}{}
\makeatother
\title[Characterizations of Lorentz Type Sobolev Multiplier Spaces ]{Characterizations of Lorentz Type Sobolev Multiplier Spaces and Their Preduals}
\author{Keng Hao Ooi}
\date{Source: arXiv:2601.12206v1 (CC BY 4.0)}
\begin{document}
\maketitle

\noindent\emph{Source and licence.} Characterizations of Lorentz Type Sobolev Multiplier Spaces and Their Preduals. Version arXiv:2601.12206v1 (CC BY 4.0), distributed under the Creative Commons Attribution 4.0 International licence, as recorded in the arXiv licence metadata for this exact version and shown on the abstract page https://arxiv.org/abs/2601.12206v1. Full rights notice: licenses/arxiv-2601.12206-CC-BY-4.0.txt.\par\smallskip

\noindent\textit{Editorial scope.} Counted unit: the Proposition whose source label is \texttt{lorentz dense} in the article (source0.tex lines 266--272, source label \texttt{lorentz dense}), delivered together with the article's own complete original proof of it (source0.tex lines 274--383, one \texttt{proof} environment of 110 lines). No mathematics is written, repaired or re-derived: statement and proof are the article's own text line for line. Required context retained uncounted: none. Every definition, display and label of those lines is retained unchanged. Omitted from this package and not redistributed: 1--265, 273--273, 384--1501. Nothing inside a retained excerpt is omitted or altered: every retained body is delivered exactly as the article's lines are written, so no omission receipt of any kind accompanies this package. Citations: the retained excerpts carry no citation command at all, so no bibliography entry of the article is transcribed and no reference is redistributed. Reference adaptation: none was needed and none was made; every reference inside the delivered unit resolves inside it, so no unresolved reference or citation is produced. Printed slips kept literal: no printed word, symbol, hypothesis or number of the counted statement or of its proof was altered. Layout and class: the article's own class is \texttt{amsart} with packages including amsmath, amstext, amssymb, amsopn, amsthm, mathrsfs, epsfig, graphics, latexsym, graphicx, color; the wrapper is the lane's pinned \texttt{amsart} editorial wrapper at 10pt on A4, which restates the theorem-like environments the retained text needs and adds the article's own macro definitions that the retained text needs (none), with extra wrapper packages (bm, bbm, dsfont, xspace, cancel, mathtools). The delivered numbering is the unit's own. Assets: no retained span contains a figure environment, a graphics-inclusion command, a PDF-page inclusion, a table or a TikZ picture; no image file is copied into this package, the package declares no asset dependency and no image file is redistributed with it. Licence: version arXiv:2601.12206v1 of this article is distributed under the Creative Commons Attribution 4.0 International licence, as recorded in the arXiv licence metadata for this exact version and shown on the abstract page https://arxiv.org/abs/2601.12206v1. A full rights notice is delivered at licenses/arxiv-2601.12206-CC-BY-4.0.txt. Characters: every retained excerpt is pure ASCII, so the delivered engine, XeLaTeX with its default Latin Modern text font, reports no missing character.
\begin{proposition}\label{lorentz dense}
Let $X$ be a locally compact Hausdorff space and $0<p,q<\infty$. Assume that $\mu$ is a $\sigma$-finite positive measure on $X$ that
\begin{align*}
X=\bigcup_{n=1}^{\infty}K_{n},
\end{align*}
where $K_{n}$ is compact with $\mu(K_{n})<\infty$, $n\in\mathbb{N}$. Denote by $L^{p,q}(\mu)$ the corresponding Lorentz space and $C_{0}$ the set of all compactly supported continuous functions on $X$. Then $C_{0}$ is dense in $L^{p,q}(\mu)$.
\end{proposition}

\begin{proof}
We first prove the set $\mathcal{K}$ of all compactly supported functions is dense in $L^{p,q}(\mu)$. Let $f\in L^{p,q}(\mu)$ and 
\begin{align*}
E_{n}=\bigcup_{i=1}^{n}K_{i},\quad n\in\mathbb{N}.
\end{align*}
For each $n\in\mathbb{N}$, the function $f\chi_{E_{n}}$ is compactly supported. Since $f\in L^{p,q}(\mu)$, we may assume that $f$ is real-valued. Besides, it is easy to argue that $\mu(\{|f|>t\})<\infty$ for each $0<t<\infty$. As a consequence,  $f-f\chi_{E_{n}}$ converges to $0$ pointwise as $n\rightarrow\infty$. The observation that 
\begin{align*}
\mu(\{|f-f\chi_{E_{n}}|>t\})\leq\mu(\{|f|>t\}),\quad 0<t<\infty,\quad n\in\mathbb{N}
\end{align*}
entails
\begin{align*}
\|f-f\chi_{E_{n}}\|_{L^{p,q}(\mu)}^{q}=p\int_{0}^{\infty}\mu(\{|f-f\chi_{E_{n}}|>t\})^{\frac{q}{p}}t^{q}\frac{dt}{t}\rightarrow 0,\quad n\rightarrow\infty
\end{align*}
by Lebesgue dominated convergence theorem. The density of $\mathcal{K}$ in $L^{p,q}(\mu)$ follows.

Subsequently, we prove that $C_{b}$ is dense in $\mathcal{K}$, where $C_{b}$ is the set of all bounded continuous functions on $X$. Let $f\in\mathcal{K}$ and $f_{M}=\max(\min(f,M),-M)$, $M>0$. By change of variables, we have
\begin{align*}
\|f_{M}-f\|_{L^{p,q}(\mu)}^{q}&=p\int_{0}^{\infty}\mu(\{|f_{M}-f|>t\})^{\frac{q}{p}}t^{q}\frac{dt}{t}\\
&=\frac{p}{q}\int_{0}^{\infty}\mu(\{|f_{M}-f|>t^{\frac{1}{q}}\})^{\frac{q}{p}}dt.
\end{align*}
Note that 
\begin{align*}
&\{|f_{M}-f|>t^{\frac{1}{q}}\}\\
&=\{\{f>M\}\cap\{|f_{M}-f|>t^{\frac{1}{q}}\}\}\cup\{\{f<-M\}\cap\{|f_{M}-f|>t^{\frac{1}{q}}\}\}\\
&=\{\{f>M\}\cap\{f-M>t^{\frac{1}{q}}\}\}\cup\{\{f<-M\}\cap\{-M-f>t^{\frac{1}{q}}\}\}\\
&=\{\{f>M\}\cap\{|f|>M+t^{\frac{1}{q}}\}\}\cup\{\{f<-M\}\cap\{|f|>M+t^{\frac{1}{q}}\}\}\\
&=\{\{|f|>M\}\cap\{|f|>M+t^{\frac{1}{q}}\}\}\\
&=\{|f|>M+t^{\frac{1}{q}}\}.
\end{align*}
Using the elementary inequality that
\begin{align*}
(a+b)^{r}\leq C(r)(a^{r}+b^{r}),\quad C(r)=\max\{2^{r-1},1\}
\end{align*}
with $r=1/q$, $a^{r}=M$, and $b=t$, one has 
\begin{align*}
\{|f|>M+t^{\frac{1}{q}}\}\subseteq\{|f|>C'(q)(M^{q}+t)^{\frac{1}{q}}\},
\end{align*}
and hence
\begin{align*}
\|f_{M}-f\|_{L^{p,q}(\mu)}^{q}&=\frac{p}{q}\int_{0}^{\infty}\mu(\{|f|>M+t^{\frac{1}{q}}\})^{\frac{q}{p}}dt\\
&\leq\frac{p}{q}\int_{0}^{\infty}\mu(\{|f|>C'(q)(M^{q}+t)^{\frac{1}{q}}\})^{\frac{q}{p}}dt\\
&=\frac{p}{q}\int_{M^{q}}^{\infty}\mu(\{|f|>C'(q)t^{\frac{1}{q}}\})^{\frac{q}{p}}dt\\
&=\frac{p}{qC''(q)}\int_{C''(q)M^{q}}^{\infty}\mu(\{|f|>t^{\frac{1}{q}}\})^{\frac{q}{p}}dt.
\end{align*}
Since
\begin{align*}
\|f\|_{L^{p,q}(\mu)}^{q}=p\int_{0}^{\infty}\mu(\{|f|>t\})^{\frac{q}{p}}t^{q}\frac{dt}{t}=\frac{p}{q}\int_{0}^{\infty}\mu(\{|f|>t^{\frac{1}{q}}\})^{\frac{q}{p}}dt<\infty,
\end{align*}
we obtain $\|f_{M}-f\|_{L^{p,q}(\mu)}^{q}\rightarrow 0$ as $M\rightarrow\infty$. For any $\varepsilon>0$, choose an $M>0$ such that $\|f_{M}-f\|_{L^{p,q}(\mu)}<\varepsilon$. Since $X$ is $\sigma$-finite, we can assume that the sequence $\{K_{n}\}_{n=1}^{\infty}$ is increasing. Pick an $n_{0}\in\mathbb{N}$ such that the interior $K_{n_{0}}^{\circ}$ of $K_{n_{0}}$ contains the support of $f$. Since $f_{M}$ is measurable and $K_{n_{0}}^{\circ}$ is of finite measure, by Lusin's theorem, there is a closed set $F\subseteq K_{n_{0}}^{\circ}$ such that $\mu(K_{n_{0}}^{\circ}\setminus F)<\varepsilon^{p}$ and $f_{M}|_{F}$ is a continuous function. Note that $f_{M}$ is identically zero on the closed set $\widetilde{F}=X\setminus K_{n_{0}}^{\circ}$. By gluing lemma, the function $f_{M}|_{F\cup\widetilde{F}}$ is continuous. Since $X$ is locally compact Hausdorff, we can choose by Tietze extension theorem a continuous function $v$ such that $|v|\leq M$ and $v=f_{M}$ on $F\cup\widetilde{F}$. Then
\begin{align*}
\|v-f_{M}\|_{L^{p,q}(\mu)}^{q}&=p\int_{0}^{\infty}\mu(\{|v-f_{M}|>t\})^{\frac{q}{p}}t^{q}\frac{dt}{t}\\
&=p\int_{0}^{\infty}\mu(\{x\in K_{n_{0}}^{\circ}\setminus F:|v(x)-f_{M}(x)|>t\})^{\frac{q}{p}}t^{q}\frac{dt}{t}\\
&=p\int_{0}^{2M}\mu(\{x\in K_{n_{0}}^{\circ}\setminus F:|v(x)-f_{M}(x)|>t\})^{\frac{q}{p}}t^{q}\frac{dt}{t}\\
&\leq C(p,q,M)\mu(K_{n_{0}}^{\circ}\setminus F)^{\frac{q}{p}}\\
&<C(p,q,M)\varepsilon^{q}.
\end{align*}
As a result, we have 
\begin{align*}
\|f-v\|_{L^{p,q}(\mu)}\leq\kappa_{0}(\|f-f_{M}\|_{L^{p,q}(\mu)}+\|v-f_{M}\|_{L^{p,q}(\mu)})<C(p,q,M,\kappa_{0})\varepsilon,
\end{align*}
where $\kappa_{0}=C(p,q)$ and the density of $C_{b}$ in $L^{p,q}(\mu)$ now follows.

Now we claim that $C_{0}$ is dense in $C_{b}$. To this end, let $v$ be continuous and $|v|\leq M$ for some $M>0$. For each $N\in\mathbb{N}$, let $O_{N}=\{|v|>1/N\}$. Then $O_{N}$ is open and 
\begin{align*}
\mu(O_{N})\leq C(p,q,N)\|f\|_{L^{p,q}(\mu)}<\infty.
\end{align*}
We observe that
\begin{align*}
\|v\chi_{O_{N}^{c}}\|_{L^{p,q}(\mu)^{q}}&=p\int_{0}^{\infty}\mu(\{x\in O_{N}^{c}:|v(x)|>t\})^{\frac{q}{p}}t^{q}\frac{dt}{t}\\
&=p\int_{0}^{\frac{1}{N}}\mu(\{x\in O_{N}^{c}:|v(x)|>t\})^{\frac{q}{p}}t^{q}\frac{dt}{t}\\
&\leq p\int_{0}^{\frac{1}{N}}\mu(\{|v|>t\})^{\frac{q}{p}}t^{q}\frac{dt}{t}\\
&\rightarrow 0
\end{align*}
as $N\rightarrow\infty$. Thus, for any $\varepsilon>0$, there is an open set $O\subseteq\mathcal{R}$ such that $\mu(O)<\infty$ and $\|v\chi_{O^{c}}\|_{L^{p,q}(\mu)}<\varepsilon$. Since $\mu$ is additive, we have 
\begin{align*}
\sum_{n=1}^{\infty}\mu(O\cap E_{n})=\mu(O)<\infty,
\end{align*}
where the sequence $\{E_{n}\}_{n=1}^{\infty}\subseteq X$ is chosen such that $\{E_{n}\}_{n=1}^{\infty}$ is disjoint and $O\cap E_{n}$ is precompact for each $n\in\mathbb{N}$. There is some $n_{0}\in\mathbb{N}$ such that
\begin{align*}
G=O\cap\left(\bigcup_{n=n_{0}}^{\infty}E_{n}\right)
\end{align*}
satisfies
\begin{align*}
\mu(G)\leq\sum_{n=n_{0}}^{\infty}\mu(O\cap E_{n})<\varepsilon^{q}.
\end{align*}
Denote by 
\begin{align*}
H=O\cap\left(\bigcup_{n=1}^{n_{0}}E_{n}\right).
\end{align*}
Since $X$ is locally compact Hausdorff and $H$ is precompact, we can choose by Urysohn's lemma an $\eta\in C_{0}$ such that $0\leq\eta\leq 1$ and $\eta=1$ on $H$. We have
\begin{align*}
&\|\eta v-v\|_{L^{p,q}(\mu)}\\
&\leq(\kappa_{0}^{2}+1)(\|(\eta v-v)\chi_{O^{c}}\|_{L^{p,q}(\mu)}+\|(\eta v-v)\chi_{G}\|_{L^{p,q}(\mu)}+\|(\eta v-v)\chi_{H}\|_{L^{p,q}(\mu)})\\
&=(\kappa_{0}^{2}+1)(\|(\eta v-v)\chi_{O^{c}}\|_{L^{p,q}(\mu)}+\|(\eta v-v)\chi_{G}\|_{L^{p,q}(\mu)}\|_{L^{p,q}(\mu)})
\end{align*}
as $\eta=1$ on $H$. On the other hand, one has
\begin{align*}
\|(\eta v-v)\chi_{O^{c}}\|_{L^{p,q}(\mu)}^{q}&=p\int_{0}^{\infty}\mu(\{x\in O^{c}:|\eta(x)v(x)-v(x)|>t\})^{\frac{q}{p}}t^{q}\frac{dt}{t}\\
&\leq p\int_{0}^{\infty}\mu\left(\left\{x\in O^{c}:|v(x)|>\frac{t}{2}\right\}\right)^{\frac{q}{p}}t^{q}\frac{dt}{t}\\
&\leq C(p,q)\|v\chi_{O^{c}}\|_{L^{p,q}(\mu)}\\
&<C(p,q)\varepsilon^{q}.
\end{align*}
Moreover, we have
\begin{align*}
\|(\eta v-v)\chi_{G}\|_{L^{p,q}(\mu)}^{q}&\leq p\int_{0}^{2M}\mu(\{x\in G:|\eta(x)v(x)-v(x)|>t\})^{\frac{q}{p}}t^{q}\frac{dt}{t}\\
&\leq C(p,q,M)\mu(G)\\
&<C(p,q,M)\varepsilon^{q},
\end{align*}
we conclude that $\|\eta v-v\|_{L^{p,q}(\mu)}<C'(p,q,M,\kappa_{0})\varepsilon$, which finishes the proof as $\eta v\in C_{0}$.
\end{proof}

\end{document}
